%% Heptagon of presentations of the chiral-topological Swiss-cheese operad.
%% Upgrade of the pentagon (five presentations) to a heptagon (seven).
%% Included in Vol II theory sequence after factorization_swiss_cheese.

%% Local macro safety (providecommand only)
\providecommand{\SCchtop}{\mathsf{SC}^{\mathrm{ch},\mathrm{top}}}
\providecommand{\Fact}{\mathrm{Fact}}
\providecommand{\CoFact}{\mathrm{CoFact}}
\providecommand{\FM}{\mathrm{FM}}
\providecommand{\Conf}{\mathrm{Conf}}
\providecommand{\Ran}{\mathrm{Ran}}
\providecommand{\Obs}{\mathsf{Obs}}
\providecommand{\Hoch}{\mathrm{Hoch}}
\providecommand{\ChirHoch}{\mathrm{ChirHoch}}
\providecommand{\Bulk}{\mathrm{Bulk}}
\providecommand{\Bdy}{\mathrm{Bdy}}
\providecommand{\MC}{\mathrm{MC}}
\providecommand{\QME}{\mathrm{QME}}
\providecommand{\End}{\mathrm{End}}
\providecommand{\Mod}{\mathrm{Mod}}
\providecommand{\id}{\mathrm{id}}
\providecommand{\cA}{\mathcal{A}}
\providecommand{\cB}{\mathcal{B}}
\providecommand{\cC}{\mathcal{C}}
\providecommand{\cE}{\mathcal{E}}
\providecommand{\cZ}{\mathcal{Z}}
\providecommand{\cM}{\mathcal{M}}
\providecommand{\fg}{\mathfrak{g}}
\providecommand{\Lie}{\mathsf{Lie}}
\providecommand{\Ass}{\mathsf{Ass}}
\providecommand{\Com}{\mathsf{Com}}
\providecommand{\QCoh}{\operatorname{QCoh}}
\providecommand{\IndCoh}{\operatorname{IndCoh}}
\providecommand{\Dmod}{D\text{-}\mathsf{mod}}
\providecommand{\colim}{\operatorname{colim}}
\providecommand{\ChirAlg}{\mathsf{ChirAlg}}
\providecommand{\Stk}{\mathsf{Stk}}
\providecommand{\RSpec}{\mathbf{RSpec}}
\providecommand{\bMap}{\mathbf{Map}}

\chapter{The Heptagon of Presentations of \texorpdfstring{$\SCchtop$}{SC}}
\label{ch:sc-chtop-heptagon}

%%%%%%%%%%%%%%%%%%%%%%%%%%%%%%%%%%%%%%%%%%%%%%%%%%%%%%%%%%%%%%%%%%%%%%%%%%%%%
\section*{Prologue}
%%%%%%%%%%%%%%%%%%%%%%%%%%%%%%%%%%%%%%%%%%%%%%%%%%%%%%%%%%%%%%%%%%%%%%%%%%%%%

One operad; seven faces. Each face is a geometric, algebraic, or
categorical presentation, and each pair of faces is connected by a named
equivalence of $\infty$-operads, $\infty$-quasi-isomorphic after
fixing a Drinfeld associator $\Phi\in\mathrm{GRT}_1(\mathbb{Q})$.
Five faces are classical: the operadic face
(Chapter~\ref{ch:factorization-swiss-cheese}), its Koszul dual
(Chapter~\ref{ch:theory-equivalence}), and the factorization,
BV/BRST, and convolution faces. Two further faces inscribe themselves:
the Drinfeld-centre face through the factorization-module category
$\Rep_{\mathrm{fact}}(\cA)$, and the derived-algebraic-geometry face
through global sections of a stack classifying $E_1$-chiral algebras
equipped with $\SCchtop$-datum.
Adjacent edges of the heptagon are named theorems; the diagonal edges
are consequences of the seven adjacencies.

The reconstitution is rigid: the $\infty$-operadic structure is
determined by any one face, and the seven faces exhaust the
presentations currently known. The absence of a conformal vector
freezes the heptagon at $\SCchtop$; its presence triggers
topologisation to $E_3$-topological (Chapter~\ref{ch:topologization}).
The heptagon is the generic object; topologisation is the exceptional
collapse.

%%%%%%%%%%%%%%%%%%%%%%%%%%%%%%%%%%%%%%%%%%%%%%%%%%%%%%%%%%%%%%%%%%%%%%%%%%%%%
\section{Setting: the object of the heptagon}
\label{sec:heptagon-setting}
%%%%%%%%%%%%%%%%%%%%%%%%%%%%%%%%%%%%%%%%%%%%%%%%%%%%%%%%%%%%%%%%%%%%%%%%%%%%%

Throughout this chapter $X$ denotes a smooth complex curve and $\cA$
denotes an $E_1$-chiral algebra on $X$. We write $Z^{\mathrm{der}}_{\mathrm{ch}}(\cA)$
for the derived chiral centre $C^{\bullet}_{\mathrm{ch}}(\cA,\cA)$
computed against the algebraic chiral endomorphism operad
$\End^{\mathrm{ch}}_\cA$ (Chapter~\ref{ch:theory-hochschild}). The pair
\begin{equation}
  \label{eq:heptagon-pair}
  \bigl(Z^{\mathrm{der}}_{\mathrm{ch}}(\cA),\;\cA\bigr)
\end{equation}
is the \emph{$\SCchtop$-datum} of $\cA$: bulk acts on boundary, the
boundary carries no action on the bulk.

\begin{convention}[What the heptagon presents]\label{conv:heptagon-object}
All seven faces present the same two-coloured coloured operad $\SCchtop$
of homotopy type
\begin{equation}
  \SCchtop(k,m)
  \;\simeq\;
  C_{\bullet}\bigl(\FM_k(\mathbb{C})\bigr)
  \,\otimes\,
  C_{\bullet}\bigl(\Conf_m(\mathbb{R})\bigr)
\end{equation}
on the open strata, with two boundary colours
(closed $=\FM_k(\mathbb{C})$, open $=\Conf_m(\mathbb{R})$) related by
codimension-one mixed stratum $\FM_k(\mathbb{C})\times\Conf_m(\mathbb{R})$,
and no open-to-closed operations. The operad is \emph{not} self-Koszul-dual
(AP-SC-NOT-SELFDUAL): $(\SCchtop)^{!}=(\Lie,\Ass,\text{shuffle-mixed})$.
\end{convention}

\begin{remark}[Directional asymmetry]\label{rem:heptagon-asymmetry}
Every face records the directional bulk-to-boundary composition;
no face admits a natural open-to-closed operation. The asymmetry forces
the Koszul dual to be non-self, forces the factorization face to act
through module categories rather than algebras, and forces the
Drinfeld-centre face to be a \emph{right adjoint} (centre), not an
averaging.
\end{remark}

%%%%%%%%%%%%%%%%%%%%%%%%%%%%%%%%%%%%%%%%%%%%%%%%%%%%%%%%%%%%%%%%%%%%%%%%%%%%%
\section{The five classical faces}
\label{sec:five-classical-faces}
%%%%%%%%%%%%%%%%%%%%%%%%%%%%%%%%%%%%%%%%%%%%%%%%%%%%%%%%%%%%%%%%%%%%%%%%%%%%%

We restate the pentagon vertices as inputs to the heptagon, fixing
notation for edges in \S\ref{sec:heptagon-edges}.

\subsection{Face (1): the operadic face}\label{subsec:face-operadic}

Face~(1) is the homotopy colored operad of
Definition~\ref{def:factorization-swiss-cheese-operad} of
Chapter~\ref{ch:factorization-swiss-cheese}: generators are the
codimension-one boundary strata of $\FM_k(\mathbb{C})\times\Conf_m(\mathbb{R})$;
relations are the codimension-two strata (Arnold for the closed
colour, Stasheff for the open colour, and the mixed Stasheff--Arnold
relation from the collision of an open point with the closed
configuration). The presentation is strict in the rational homotopy
category after fixing $\Phi$ (Kontsevich--Tamarkin formality).

\subsection{Face (2): the Koszul-dual face}\label{subsec:face-koszul-dual}

Face~(2) is the $\infty$-operadic cofibrant replacement
\begin{equation}
  W(\SCchtop)
  \;:=\;
  \Omega\bigl(B(\SCchtop)\bigr)
\end{equation}
via Ginzburg--Kapranov coloured bar--cobar, after fixing a Drinfeld
associator. Its underlying graded coloured operad is
$(\SCchtop)^{!}=(\Lie,\Ass,\text{shuffle-mixed})$
(Proposition~\ref{prop:sc-koszul-dual-three-sectors}).
The closed sector is
$E_2^{!}\simeq E_2\{1\}$ (Getzler--Jones self-duality up to shift),
which on the associated graded collapses to $\Lie$.
The open sector is self-dual $\Ass^{!}=\Ass$. The mixed sector is
shuffle-mixed with empty open-to-closed operations; the
closed-to-open is inherited from $\SCchtop$'s directionality.

\subsection{Face (3): the factorization face}\label{subsec:face-factorization}

Face~(3) is the factorization-algebra model. For each open $U\subset X$
the chiral observables
\begin{equation}
  \Obs^{\mathrm{ch}}_\cA(U),\qquad U\subset X,
\end{equation}
assemble into a factorization algebra on $\Ran(X)$; its derived
centre $Z^{\mathrm{der}}_{\mathrm{ch}}(\cA)$ acts on $\cA$ through the
universal brace of Chapter~\ref{ch:theory-hochschild}. The pair in
\eqref{eq:heptagon-pair} is the factorization-algebraic
$\SCchtop$-datum. This face is native to Beilinson--Drinfeld
(algebraic) and Costello--Gwilliam (differential-geometric); the
equivalence is Proposition~\ref{prop:bd-cg-equivalence} of
Chapter~\ref{ch:factorization-swiss-cheese}.

\subsection{Face (4): the BV/BRST face}\label{subsec:face-bv-brst}

Face~(4) is the observables of a bulk--boundary BV theory on
$X\times\mathbb{R}_{\geq 0}$ with boundary condition $\cA$. The bulk
action $S_{\mathrm{bulk}}$ satisfies the classical master equation;
after BV quantisation, the quantum observables on the closed half
(two-disc topology) carry the logarithmic
$\SCchtop$-algebra structure of
Definition~\ref{def:log-SC-algebra}, the open sector records
boundary local operators, and the mixed sector is line-operator
insertion of bulk into boundary. The QME on the combined bulk--boundary
system is the open/closed Maurer--Cartan equation associated to
$\SCchtop$. The $\SCchtop$-datum is the \emph{pair}, not
$\cA$ alone.

\subsection{Face (5): the convolution face}\label{subsec:face-convolution}

Face~(5) is the $L_\infty$-algebra of deformations of a twisting
morphism from $B(\SCchtop)$ to $\End^{\mathrm{ch}}$. Fix a cofibrant
replacement $\cE\twoheadrightarrow\SCchtop$; the convolution
$L_\infty$-algebra
\begin{equation}
  \fg^{\,\SCchtop}_{\cA,Z}
  \;:=\;
  \mathrm{Hom}^{\mathrm{col}}\bigl(B(\SCchtop),\,\End^{\mathrm{ch}}_{(\cA,Z)}\bigr)
\end{equation}
controls colored $\SCchtop$-algebra structures on the pair
$(\cA,Z^{\mathrm{der}}_{\mathrm{ch}}(\cA))$ through its Maurer--Cartan set
(Vallette, Loday--Vallette for the colored case). The
$\SCchtop$-datum is a Maurer--Cartan element $\alpha_T\in
\fg^{\,\SCchtop}_{\cA,Z}$; gauge equivalence is $\Phi$-dependent;
the gauge class is $\Phi$-independent and equals the operadic
face's homotopy class.

%%%%%%%%%%%%%%%%%%%%%%%%%%%%%%%%%%%%%%%%%%%%%%%%%%%%%%%%%%%%%%%%%%%%%%%%%%%%%
\section{Face (6): Drinfeld centre of the factorization module category}
\label{sec:face-drinfeld-centre}
%%%%%%%%%%%%%%%%%%%%%%%%%%%%%%%%%%%%%%%%%%%%%%%%%%%%%%%%%%%%%%%%%%%%%%%%%%%%%

The sixth face promotes the factorization face to its categorical
shadow. In the categorical direction, the $\SCchtop$-datum is
equivalent to the $E_2$-monoidal structure on a module category for
$\cA$, together with a half-braiding whose spectral reduction is
Face~(3)'s $r$-matrix. The construction is a right adjoint, not an
averaging; the half-braiding is
\emph{constructed}, not narrated.

\subsection{The factorization module category}\label{subsec:fact-mod-cat}

\begin{definition}[Factorization module category]\label{def:rep-fact}%
\index{factorization module category!$\Rep_{\mathrm{fact}}$|textbf}%
\index{$\Rep_{\mathrm{fact}}$}
Let $\cA$ be an $E_1$-chiral algebra on $X$. The \emph{factorization
module category}
\begin{equation}
  \Rep_{\mathrm{fact}}(\cA)
  \;:=\;
  \mathbb{A}\text{-}\mathrm{fact}\mathrm{Mod}(\Dmod(\Ran(X)))
\end{equation}
is the dg category of factorization modules for $\cA$: objects are
$\cA$-module $D$-modules on $\Ran(X)$ whose factorization structure
is compatible with the chiral $\cA$-module action. Morphisms are
$\cA$-linear factorization morphisms. Composition is defined
using the chiral tensor product
$\otimes^{\mathrm{ch}}$ on $\Ran(X)$
(Beilinson--Drinfeld, GR17 IV.5; Francis--Gaitsgory 2012).
\end{definition}

\begin{remark}[Three local models agree]\label{rem:three-local-models}
On the formal disc $D_x\subset X$, the category
$\Rep_{\mathrm{fact}}(\cA)$ restricts to the usual category
$\cA_{x}\text{-}\Mod$ of modules over the $\cA$-mode algebra
(the Lie algebra of modes of $\cA$ on $D_x$). On the punctured
formal disc $D^\times_x$ it restricts to the category of smooth
$\hat{\cA}_{x}$-modules. The three models glue along
$D^\times_x\hookrightarrow D_x$ via the pullback of the factorization
structure; the glue is the content of
Lemma~\ref{lem:rep-fact-local-global}.
\end{remark}

\begin{lemma}[Local--global description of $\Rep_{\mathrm{fact}}$;
\ClaimStatusProvedElsewhere]\label{lem:rep-fact-local-global}
$\Rep_{\mathrm{fact}}(\cA)$ is equivalent to the (∞,1)-limit
of the factorization presheaf of local module categories over the
Ran space $\Ran(X)$. Equivalently, it is the category of
modules over the factorization algebra $\Obs^{\mathrm{ch}}_\cA$ in
$\Dmod(\Ran(X))$. The equivalence is Beilinson--Drinfeld for the
algebraic incarnation, Costello--Gwilliam for the analytic one; see
Gaitsgory--Rozenblyum~\cite{GR17} IV.5.
\end{lemma}

\begin{proof}[Attribution]
Beilinson--Drinfeld \textit{Chiral Algebras} Ch.~3, Lemma~3.4.15;
Gaitsgory--Rozenblyum, \textit{Study in Derived Algebraic Geometry},
Vol.~II, Ch.~IV.5, for the $\infty$-categorical form. The
factorization-module description is in Francis--Gaitsgory 2012
Definition~2.7.6.
\end{proof}

\subsection{The half-braiding construction}\label{subsec:half-braiding}

\begin{definition}[Chiral half-braiding]\label{def:chiral-half-braiding}%
\index{half-braiding!chiral|textbf}%
\index{factorization module category!half-braiding}
Let $\cA$ be an $E_1$-chiral bialgebra with spectral coproduct
$\Delta_z\colon \cA\to\cA\otimes\cA$
(Chapter~\ref{ch:theory-hochschild}, Definition~\ref{def:half-braiding-explicit}).
For $N\in\Rep_{\mathrm{fact}}(\cA)$ and $a\in\cA$, $n\in N$, set
\begin{equation}
  \label{eq:chiral-half-braiding}
  \sigma_\cA(z)(a\otimes n)
  \;:=\;
  \sum\bigl(\Delta_z(a)\bigr)_{(2)}\cdot n
  \;\otimes\;
  \bigl(\Delta_z(a)\bigr)_{(1)}
  \;=\;
  \sum a_{(2),z}\cdot n\otimes a_{(1),z}.
\end{equation}
The family $\{\sigma_\cA(z)\}_{z\in X\setminus\{0\}}$ assembles into a
natural isomorphism of functors
$\cA\otimes(-)\Rightarrow(-)\otimes\cA$ on
$\Rep_{\mathrm{fact}}(\cA)$, meromorphic in the spectral parameter $z$,
with first-order pole at $z=0$.
\end{definition}

\begin{proposition}[Half-braiding from $\Delta_z$; \ClaimStatusProvedHere]
\label{prop:half-braiding-construction}%
\index{half-braiding!from $\Delta_z$}
The half-braiding $\sigma_\cA(z)$ of
Definition~\ref{def:chiral-half-braiding} satisfies:
\begin{enumerate}[label=\textup{(\roman*)}]
  \item hexagon identities (quasi-triangularity of $\cA$ expressed
  on factorization modules);
  \item spectral $R$-matrix:
  $\sigma_\cA(z)|_{V_u,V_v}=R(u-v)$ on evaluation modules
  $V_u,V_v$ with spectral parameters $u,v$;
  \item $\Rep_{\mathrm{fact}}(\cA)$ becomes a braided monoidal
  $\infty$-category, i.e.\ an $E_2$-algebra in $\mathrm{Cat}_\infty$
  (MP1/FM41); $\cA$ itself remains an $E_1$-chiral algebra;
  \item the braiding \emph{extends} to all of $\Rep_{\mathrm{fact}}(\cA)$,
  not only to evaluation modules: for generic $M,N$ in
  $\Rep_{\mathrm{fact}}$, $\sigma_\cA$ provides a meromorphic
  natural transformation $M\otimes N\to N\otimes M$ with poles
  on the discriminant of the module supports.
\end{enumerate}
\end{proposition}

\begin{proof}
(i) follows from the defining quasi-triangularity axiom for
$\Delta_z$ (Chapter~\ref{ch:theory-drinfeld}): for
$a,b\in\cA$,
\(
  \sigma_\cA(z)(ab\otimes n)=
  \sigma_\cA(z)(a\otimes\sigma_\cA(z)(b\otimes n))
\)
after applying the spectral coproduct to the product $ab$ and reading
off Sweedler indices. The axiomatics match the standard quasi-Hopf
hexagon after translation to the half-braiding form.

(ii) is the evaluation specialisation of~\eqref{eq:chiral-half-braiding}:
on $V_u\otimes V_v$, $\Delta_z(a)\cdot(v_u\otimes v_v)$ evaluates to
$\pi_u\otimes\pi_v(\Delta_z(a))$, which is the spectral
$R$-matrix acting on the tensor product modulo the swap, recovering
$R(u-v)$ by the Drinfeld--Jimbo reconstruction of $R$ from
$\Delta_z$ (Chapter~\ref{ch:ordered-chiral}).

(iii) is a consequence of (i) + (ii) + the module category structure
of $\Rep_{\mathrm{fact}}$: a braided monoidal structure is a
half-braiding that is itself natural, and naturality is the
factorization-module compatibility of Lemma~\ref{lem:rep-fact-local-global}.
By Lurie HA~5.1 the braided monoidal structure is an
$E_2$-algebra in $\mathrm{Cat}_\infty$.

(iv) is the extension from evaluation modules to arbitrary
factorization modules, enabled by the universal property of
$\Rep_{\mathrm{fact}}(\cA)$ (Lemma~\ref{lem:rep-fact-local-global}):
every module is a colimit of evaluation modules along the factorization
structure, and $\sigma_\cA(z)$ is natural with respect to factorization
morphisms. Meromorphicity of $\sigma_\cA(z)$ on the
supports follows from OPE pole bounds on $\Delta_z$.
\end{proof}

\subsection{Drinfeld centre and factorization centre}\label{subsec:dcz-vs-fcz}

\begin{definition}[Drinfeld centre of $\Rep_{\mathrm{fact}}(\cA)$]
\label{def:drinfeld-centre-rep-fact}%
\index{Drinfeld centre!of $\Rep_{\mathrm{fact}}$|textbf}
The \emph{Drinfeld centre}
$\cZ\bigl(\Rep_{\mathrm{fact}}(\cA)\bigr)$ is the $\infty$-category of
pairs $(N,\sigma_N)$, where $N\in\Rep_{\mathrm{fact}}(\cA)$ and
$\sigma_N\colon N\otimes(-)\Rightarrow(-)\otimes N$ is a half-braiding
of $N$ compatible with $\sigma_\cA(z)$. Composition is component-wise.
The Drinfeld centre carries a canonical $E_2$-monoidal structure.
\end{definition}

\begin{theorem}[Drinfeld centre face of $\SCchtop$;
\ClaimStatusProvedHere]
\label{thm:drinfeld-centre-sc-face}%
\index{SC!Drinfeld centre face|textbf}%
\index{heptagon!face 6}
Let $\cA$ be an $E_1$-chiral bialgebra on $X$ with conformal-vector
hypothesis dropped (generic case). Then
\begin{equation}
  \label{eq:drinfeld-centre-face}
  \cZ\bigl(\Rep_{\mathrm{fact}}(\cA)\bigr)
  \;\simeq\;
  \bigl(\Rep_{\mathrm{fact}}(Z^{\mathrm{der}}_{\mathrm{ch}}(\cA))\bigr)^{E_2}
\end{equation}
as braided monoidal $\infty$-categories, where the left side is the
categorical Drinfeld centre and the right side is the category of
factorization modules over the derived chiral centre
(Definition~\ref{def:rep-fact} applied to the bulk algebra), with
its induced $E_2$-monoidal structure. The equivalence carries the
half-braiding of~\eqref{eq:chiral-half-braiding} on the left to the
spectral braiding of the right induced from
$\Delta_z^{\,Z}\colon Z^{\mathrm{der}}_{\mathrm{ch}}(\cA)\to
Z^{\mathrm{der}}_{\mathrm{ch}}(\cA)\otimes Z^{\mathrm{der}}_{\mathrm{ch}}(\cA)$.
The pair
\(
  \bigl(\cZ(\Rep_{\mathrm{fact}}(\cA)),\Rep_{\mathrm{fact}}(\cA)\bigr)
\)
is a categorified $\SCchtop$-datum: the bulk category acts on
the boundary category, with no open-to-closed direction.
\end{theorem}

\begin{proof}
Construct the equivalence in four steps.

\emph{Step 1: from half-braiding to central object.} The universal
half-braiding~\eqref{eq:chiral-half-braiding} exhibits $\cA$ itself as
\emph{not} an object of $\cZ(\Rep_{\mathrm{fact}}(\cA))$ (MP1/FM41);
instead, every element of the bar complex $B^{\,\mathrm{ch}}(\cA)$,
twisted by the $R$-matrix to compensate ordering ambiguity, lifts to a
central object. By direct computation, the lift of the vacuum gives
the identity central object; the lifts of generators of $\cA$ give
the half-braiding data of Proposition~\ref{prop:half-braiding-construction}.

\emph{Step 2: the categorified bar--cobar adjunction.} The adjoint pair
\begin{equation}
  \mathrm{forget}\colon\cZ(\Rep_{\mathrm{fact}}(\cA))\rightleftarrows
  \Rep_{\mathrm{fact}}(\cA)\colon\mathrm{centre}
\end{equation}
is the right-adjoint-to-forgetful construction (Majid 1998;
Lurie HA~5.3): the centre is the categorified \emph{commutant}
, not an averaging. Its right adjoint produces, for any
factorization module $N$, the universal half-braided lift
$\cZ(N)\in\cZ(\Rep_{\mathrm{fact}}(\cA))$.

\emph{Step 3: identification with bulk modules.} The derived chiral
centre $Z^{\mathrm{der}}_{\mathrm{ch}}(\cA)=C^{\bullet}_{\mathrm{ch}}(\cA,\cA)$
acts on $\cA$ through its $E_2$-chiral structure (chiral Higher
Deligne, chiral-Hochschild incarnation). By
Beilinson--Drinfeld--Zhu--Francis--Gaitsgory--Rozenblyum factorization
module theory (Lurie HA~5.3.1.30; BZFN10),
$\Rep_{\mathrm{fact}}(Z^{\mathrm{der}}_{\mathrm{ch}}(\cA))$ acquires the
induced braided monoidal structure from
$\Delta_z^{\,Z}$. Step~2's right adjoint then factors through this
category: $\cZ(\Rep_{\mathrm{fact}}(\cA))\simeq
\Rep_{\mathrm{fact}}(Z^{\mathrm{der}}_{\mathrm{ch}}(\cA))$.

\emph{Step 4: matching the half-braiding to the spectral braiding.}
On evaluation modules the matching is
Proposition~\ref{prop:half-braiding-construction}~(ii). The
extension to all factorization modules is
Proposition~\ref{prop:half-braiding-construction}~(iv). The two
braidings agree as natural transformations
by the universal property of factorization module categories
(Lemma~\ref{lem:rep-fact-local-global}): both compute the action of
the central object on a general module, through the
factorization-compatible chiral tensor product.

This completes the equivalence in~\eqref{eq:drinfeld-centre-face}.
The $\SCchtop$-datum at the categorical level is read off directly:
the bulk category is $\cZ(\Rep_{\mathrm{fact}}(\cA))$, the boundary
category is $\Rep_{\mathrm{fact}}(\cA)$, the module action of the
first on the second is the centre-to-forgetful counit, no open-to-closed
exists (the centre receives no extra input from the boundary beyond
what the forgetful sends).
\end{proof}

\begin{remark}[Associator dependence]\label{rem:face-6-associator}
Theorem~\ref{thm:drinfeld-centre-sc-face} is
$\infty$-quasi-isomorphic with the factorization face after fixing a
Drinfeld associator $\Phi\in\mathrm{GRT}_1(\mathbb{Q})$. The underlying
braided monoidal $\infty$-category is $\Phi$-independent; the strict
chain-level presentation depends on $\Phi$. On the Koszul locus, two
associators yield strictly equivalent half-braidings after applying
the $\Phi_1\Phi_2^{-1}$ twist.
\end{remark}

%%%%%%%%%%%%%%%%%%%%%%%%%%%%%%%%%%%%%%%%%%%%%%%%%%%%%%%%%%%%%%%%%%%%%%%%%%%%%
\section{Face (7): derived algebraic geometry}
\label{sec:face-dag}
%%%%%%%%%%%%%%%%%%%%%%%%%%%%%%%%%%%%%%%%%%%%%%%%%%%%%%%%%%%%%%%%%%%%%%%%%%%%%

The seventh face reads $\SCchtop$ as the algebra of global sections
of a derived stack classifying $E_1$-chiral algebras with
$\SCchtop$-datum. The stack is $1$-shifted symplectic in the sense
of Pantev--To\"en--Vaqui\'e--Vezzosi (PTVV~2013), and the convolution
face is recovered as the tangent Lie algebra.

\subsection{The classifying stack}\label{subsec:classifying-stack}

\begin{definition}[Stack of $E_1$-chiral algebras with
$\SCchtop$-datum]\label{def:moduli-sc-ch-top}%
\index{moduli stack!of $E_1$-chiral algebras with $\SCchtop$-datum|textbf}%
\index{$\cM_{\SCchtop}$}
Let $X$ be a smooth complex curve. Define the functor
\begin{equation}
  \cM_{\SCchtop,X}\colon
  (\mathrm{dgArt}/\mathbb{C})^{\mathrm{op}}
  \;\longrightarrow\;
  \mathrm{sSet}
\end{equation}
from derived Artinian $\mathbb{C}$-algebras to simplicial sets by
sending $R\in\mathrm{dgArt}$ to the mapping space
\begin{equation}
  \cM_{\SCchtop,X}(R)
  \;:=\;
  \bMap_{\mathrm{OperadCat}_\infty}
  \bigl(\SCchtop,\,
  \mathrm{End}^{\mathrm{col}}_{(\cA,\,Z)}\otimes_{\mathbb{C}} R\bigr),
\end{equation}
where the mapping space is taken in coloured $\infty$-operads on
$\Dmod(\Ran(X))$ and the colimit is over all pairs
$(\cA,Z)$ with $\cA$ an $R$-flat $E_1$-chiral algebra on $X$ and
$Z$ an $E_2$-chiral algebra acting on $\cA$ through the brace
structure. The stack is the sheafification of this functor with respect
to the étale topology on $\mathrm{dgArt}$.
\end{definition}

\begin{lemma}[Representability; \ClaimStatusProvedElsewhere]
\label{lem:moduli-sc-representable}
The functor $\cM_{\SCchtop,X}$ is a locally finite presented derived
Artin stack over $\mathbb{C}$. It carries a natural action of the
Drinfeld--Grothendieck--Teichm\"uller group
$\mathrm{GRT}_1(\mathbb{Q})$ through associator transport.
\end{lemma}

\begin{proof}[Attribution]
Representability: To\"en--Vaqui\'e (Compos.\ Math.\ 143, 2007,
\S3.3) for operadic moduli; Lurie, \emph{Derived Algebraic Geometry IV}
for the Artin criteria; Francis--Gaitsgory 2012 for the
factorization variant. $\mathrm{GRT}_1$-action via Willwacher 2015
(Invent.\ Math.\ 200).
\end{proof}

\begin{remark}[Three strata]\label{rem:moduli-strata}
$\cM_{\SCchtop,X}$ stratifies by conformal-vector presence:
the open substack $\cM^{\mathrm{conf}}_{\SCchtop,X}$ where the
boundary $\cA$ carries a conformal vector at non-critical level;
the closed complement $\cM^{\mathrm{nc}}_{\SCchtop,X}$ without
conformal vector; and the stratum of critical level
$\cM^{\mathrm{crit}}_{\SCchtop,X}$ where the conformal vector
degenerates. On $\cM^{\mathrm{conf}}$ the structure topologises to
$E_3$-topological (Chapter~\ref{ch:topologization}); on
$\cM^{\mathrm{nc}}$ and $\cM^{\mathrm{crit}}$ the heptagon remains
intact.
\end{remark}

\subsection{Shifted symplectic structure}\label{subsec:shifted-symplectic}

\begin{theorem}[$1$-shifted symplectic structure on
$\cM_{\SCchtop,X}$; \ClaimStatusProvedHere]
\label{thm:moduli-sc-shifted-symplectic}%
\index{PTVV!on moduli of $\SCchtop$-data}%
\index{$1$-shifted symplectic!on $\cM_{\SCchtop,X}$}
The derived stack $\cM_{\SCchtop,X}$ carries a natural
$1$-shifted symplectic form $\omega_{\SCchtop}\in
\mathrm{Sp}^1(\cM_{\SCchtop,X})$, constructed as the pairing
\begin{equation}
  \label{eq:moduli-shifted-sympl}
  \omega_{\SCchtop}
  \;=\;
  \int_{X\times\mathbb{R}_{\geq 0}}\;
  \mathrm{ev}^{*}\,\kappa_{\SCchtop},
\end{equation}
where $\mathrm{ev}$ is the evaluation map from
$\cM_{\SCchtop,X}\times X\times\mathbb{R}_{\geq 0}$ to the
universal $E_3$-chiral target (the tangent pairing of the universal
$\SCchtop$-action), and $\kappa_{\SCchtop}$ is the coloured bilinear
form on $(\SCchtop)^{!}=(\Lie,\Ass,\text{shuffle-mixed})$ given by
the Killing form on the closed $\Lie$ colour and the trace form on
the open $\Ass$ colour. The pairing is anti-symmetric of degree $1$,
and closed as a shifted $2$-form by PTVV~\cite{PTVV13}
Theorem~2.5 (mapping-stack version) applied to the source
$X\times\mathbb{R}_{\geq 0}$, which is a $2$-orientable $\infty$-pair
of dimensions $(1,1)$.
\end{theorem}

\begin{proof}
Apply PTVV~\cite{PTVV13} Theorem~2.5 to the mapping stack
$\cM_{\SCchtop,X}=\bMap(X\times\mathbb{R}_{\geq 0},
B\SCchtop{\text{-}}\mathrm{Alg})$, where
$B\SCchtop\mathrm{-Alg}$ denotes the classifying $(\infty,2)$-stack of
$\SCchtop$-algebras in $\Dmod(\Ran(X))$. The source
$X\times\mathbb{R}_{\geq 0}$ is a smooth $\infty$-pair
$(2,1)$-dimensional; the target carries a $2$-shifted
symplectic structure induced by the coloured bilinear form
$\kappa_{\SCchtop}$. Integration over the source drops degree by
$2-1=1$. Closedness follows from the closed nature of the source-side
$2$-form and PTVV's transgression machinery; non-degeneracy
follows from the non-degeneracy of
$\kappa_{\SCchtop}$ (the $\Lie$-sector Killing form is non-degenerate
on semisimple $\Lie$-sectors; the $\Ass$-trace is non-degenerate on
finite-dimensional $\Ass$-sectors; the mixed-sector shuffle form
inherits non-degeneracy from the tensor product of the two).
\end{proof}

\begin{remark}[Why the pair (2,1)]\label{rem:why-pair-21}
The source $X\times\mathbb{R}_{\geq 0}$ carries two distinct
dimensions: $X$ is $1$-complex-dimensional hence $2$-real-dimensional,
$\mathbb{R}_{\geq 0}$ is $1$-real-dimensional with boundary. PTVV's
$\infty$-pair framework allows this mixed dimension, and the
integration pairs a $1$-form on the source with a $2$-form on the
target, yielding a $1$-shifted output.
\end{remark}

\subsection{Lagrangian correspondence to the convolution face}\label{subsec:lagrangian-corr}

\begin{theorem}[Convolution face as Lagrangian section;
\ClaimStatusProvedHere]
\label{thm:convolution-lagrangian}%
\index{convolution!Lagrangian section of $\cM_{\SCchtop,X}$}%
\index{heptagon!edge (5)--(7)}
The convolution $L_\infty$-algebra
$\fg^{\,\SCchtop}_{\cA,Z}$ of Face~(5) is the tangent dg-Lie
algebra to $\cM_{\SCchtop,X}$ at the point represented by
the pair $(\cA,Z^{\mathrm{der}}_{\mathrm{ch}}(\cA))$. The Maurer--Cartan
locus
\begin{equation}
  \MC\bigl(\fg^{\,\SCchtop}_{\cA,Z}\bigr)
  \;\subset\;
  \cM_{\SCchtop,X}
\end{equation}
is a Lagrangian subvariety of
$(\cM_{\SCchtop,X},\omega_{\SCchtop})$ in the PTVV sense: the
pullback of $\omega_{\SCchtop}$ to the MC locus vanishes and the
normal bundle of MC inside $\cM_{\SCchtop,X}$ is anti-symplectic to
the tangent of MC.
\end{theorem}

\begin{proof}
Tangent identification: by the universal property of mapping stacks,
the tangent complex of $\cM_{\SCchtop,X}$ at $(\cA,Z)$ is
\begin{equation}
  T_{(\cA,Z)}\cM_{\SCchtop,X}
  \;=\;
  \mathrm{Hom}^{\mathrm{col}}\bigl(B(\SCchtop),\,\End^{\mathrm{col}}_{(\cA,Z)}\bigr),
\end{equation}
which is exactly $\fg^{\,\SCchtop}_{\cA,Z}$ (Face~(5); Loday--Vallette
\textsection 10.2 for the coloured convolution adjoint). The
$L_\infty$-structure on the tangent complex matches the
convolution structure by Kontsevich formality of $\SCchtop$.

Lagrangian property: the anchored form
$\omega_{\SCchtop}|_{\MC}$ vanishes because the MC equation in
$\fg^{\,\SCchtop}_{\cA,Z}$ is equivalent to the first-order
vanishing of $\omega_{\SCchtop}$ along the tangent direction (PTVV~2013
Proposition~2.8). Anti-symplectic normal: the normal bundle is
$\fg^{\,\SCchtop}_{\cA,Z}[1]$, shifted by one, and inherits the
anti-symplectic pairing from the PTVV cotangent adjunction.
\end{proof}

\begin{corollary}[Face~(7) presents $\SCchtop$;
\ClaimStatusProvedHere]
\label{cor:face-7-presents-sc}
The $\SCchtop$-datum associated to any point $(\cA,Z)\in
\cM_{\SCchtop,X}$ is fully determined by the germ of
$\omega_{\SCchtop}$ at that point together with the Lagrangian
correspondence of Theorem~\ref{thm:convolution-lagrangian}.
Hence $\cM_{\SCchtop,X}$ presents $\SCchtop$ as a shifted
symplectic geometry, and recovers the operadic face through its
tangent complex.
\end{corollary}

\begin{proof}
Combine Theorem~\ref{thm:moduli-sc-shifted-symplectic} with
Theorem~\ref{thm:convolution-lagrangian}: the operadic structure of
$\SCchtop$ is encoded in the shifted symplectic form
$\omega_{\SCchtop}$; the $L_\infty$-structure is encoded in the
tangent complex; the $\SCchtop$-data as operad-algebras are in
bijection with Lagrangian sections of $\omega_{\SCchtop}$.
\end{proof}

%%%%%%%%%%%%%%%%%%%%%%%%%%%%%%%%%%%%%%%%%%%%%%%%%%%%%%%%%%%%%%%%%%%%%%%%%%%%%
\section{The seven edges of the heptagon}
\label{sec:heptagon-edges}
%%%%%%%%%%%%%%%%%%%%%%%%%%%%%%%%%%%%%%%%%%%%%%%%%%%%%%%%%%%%%%%%%%%%%%%%%%%%%

Number the faces $1,\dots,7$ as above. We prove the seven adjacent
edges as named theorems; the remaining $\binom{7}{2}-7=14$ diagonal
identifications follow by composition, hence are recorded as
remarks.

\begin{convention}[Edge regime]\label{conv:edge-regime}
Every edge is an $\infty$-quasi-isomorphism of coloured
$\infty$-operads after fixing a Drinfeld associator
$\Phi\in\mathrm{GRT}_1(\mathbb{Q})$. Strict chain-level equivalences
hold on the Koszul locus; away from it, the equivalences are
$\infty$-quasi-isomorphisms that become strict after passing to
$\Phi$-transported presentations. We suppress $\Phi$ from edge names;
each name carries its $\Phi$-dependence implicitly.
\end{convention}

\subsection{Edge (1)--(2): operadic--Koszul-dual}\label{subsec:edge-12}

\begin{theorem}[Ginzburg--Kapranov Koszul duality for
$\SCchtop$; \ClaimStatusProvedHere]\label{thm:edge-12}%
\index{heptagon!edge (1)--(2)}
The bar--cobar adjunction
$(\Omega,B)$ on coloured $\infty$-operads restricts to a
Quillen equivalence between Face~(1) and Face~(2):
\begin{equation}
  W(\SCchtop)=\Omega(B(\SCchtop))
  \;\xrightarrow{\sim}\;
  \SCchtop.
\end{equation}
The bar coalgebra $B(\SCchtop)$ is explicitly computed (closed
sector $B(\Com)=\Lie^c$, open sector $B(\Ass)=\Ass^c$, mixed sector
shuffle-bar); cobarring recovers $\SCchtop$.
\end{theorem}

\begin{proof}
The result is coloured Ginzburg--Kapranov: Vallette
\textit{A Koszul duality for properads} (Trans.\ AMS 359, 2007)
Theorem~4.1.3 for the coloured case; Hoefel~2009 (arXiv:0809.4623)
and Hoefel--Livernet~2012 (arXiv:1207.2307) specialise to the Swiss
cheese operad and verify Koszulity, producing the coloured Koszul
dual as $(\Lie,\Ass,\text{shuffle-mixed})$. The proof is
Proposition~\ref{prop:sc-koszul-dual-three-sectors} of
Chapter~\ref{ch:theory-equivalence}. The bar and
cobar functors preserve weak equivalences of $\infty$-operads
(Fresse, \textit{Modules over Operads and Functors}, Springer 2009),
hence the Quillen equivalence.
\end{proof}

\subsection{Edge (2)--(3): Koszul-dual--factorization}\label{subsec:edge-23}

\begin{theorem}[Chiral Deligne--Tamarkin for $\SCchtop$ at
$E_2$-cohomology; \ClaimStatusProvedHere]\label{thm:edge-23}%
\index{heptagon!edge (2)--(3)}
The cofibrant replacement $W(\SCchtop)$ acts on the pair
$(Z^{\mathrm{der}}_{\mathrm{ch}}(\cA),\cA)$ via the brace structure of
Chapter~\ref{ch:theory-hochschild}, and this action is
$\infty$-quasi-isomorphic to the factorization action of
Face~(3). Explicitly,
\begin{equation}
  W(\SCchtop)
  \;\xrightarrow{\sim}\;
  \End^{\mathrm{col}}_{(\cA,Z^{\mathrm{der}}_{\mathrm{ch}}(\cA))},
\end{equation}
with the $\infty$-quasi-isomorphism fixed by $\Phi\in\mathrm{GRT}_1$.
At $E_2$-cohomology (i.e.\ passing to derived brace structure) the
equivalence is $\Phi$-independent.
\end{theorem}

\begin{proof}
Combine:
(i) the chiral higher Deligne theorem
(Theorem~\ref{thm:chiral-higher-deligne}, Chapter~\ref{ch:theory-hochschild}):
$Z^{\mathrm{der}}_{\mathrm{ch}}(\cA)$ carries a canonical
$E_2$-chiral structure;
(ii) the action of $W(\SCchtop)$ on the bulk--boundary pair by
colored brace operations (Tamarkin 2000 for the uncoloured
$E_2$; chiral upgrade via factorization brace);
(iii) Koszul-dual presentation: $W(\SCchtop)=\Omega(B(\SCchtop))$
via Edge~(1)--(2); the action factors through the cofibrant
replacement.

$\Phi$-independence at cohomology: follows from the
GRT$_1$-action on brace operations preserving the cohomology class
(Willwacher 2015 for uncoloured $E_2$; coloured analogue via
coloured Kontsevich formality).
\end{proof}

\subsection{Edge (3)--(4): factorization--BV/BRST}\label{subsec:edge-34}

\begin{theorem}[Observables of BV theory present the factorization
face; \ClaimStatusProvedHere]\label{thm:edge-34}%
\index{heptagon!edge (3)--(4)}
For a bulk--boundary BV theory on $X\times\mathbb{R}_{\geq 0}$ with
boundary condition $\cA$, Costello--Gwilliam observables
$\Obs^{\mathrm{q}}(\cdot)$ assemble into a factorization algebra
on $X\times\mathbb{R}_{\geq 0}$ whose restriction to the closed half
is $\Obs^{\mathrm{ch}}_\cA$ of Face~(3). The QME for the bulk--boundary
action is the open/closed Maurer--Cartan equation for the pair
$(\cA,Z^{\mathrm{der}}_{\mathrm{ch}}(\cA))$.
\end{theorem}

\begin{proof}
Costello--Gwilliam \textit{Factorization Algebras in QFT}~Vol.~II
Theorem~11.3.3 (closed sector); the bulk--boundary extension is
Butson--Yoo 2018 (arXiv:1807.02341); the QME/MC identification is
the logarithmic $\SCchtop$-algebra structure of
Definition~\ref{def:log-SC-algebra} (Chapter~\ref{ch:bv-construction}):
the QME reads
$\partial S_{\mathrm{bulk-bdy}} + \tfrac{1}{2}\{S_{\mathrm{bulk-bdy}},
S_{\mathrm{bulk-bdy}}\}=0$ with coloured BV bracket extending the bulk
BV bracket by a boundary-to-bulk bracket, and this is the MC equation
for the pair, not for $\cA$ alone.
\end{proof}

\begin{proposition}[Heptagon edge $(3)\leftrightarrow(4)$, propositional
form; \ClaimStatusProvedHere]\label{prop:heptagon-edge-34}%
\index{heptagon!edge (3)--(4) propositional form}%
The factorization presentation (Face~(3)) and the BV/BRST presentation
(Face~(4)) of $\SCchtop$ determine the same $\infty$-operadic data on
$(\cA, Z^{\mathrm{der}}_{\mathrm{ch}}(\cA))$. Equivalently, the
$L_\infty$-convolution algebra $\fg^{\,\SCchtop}_{\cA,Z}$ of Face~(5)
is quasi-isomorphic, through the BV/BRST presentation of Face~(4), to
the factorization-homology universal brace action of Face~(3); the
quasi-isomorphism is implemented by the Swiss-cheese formality of
Kontsevich--Tamarkin after fixing a Drinfeld associator $\Phi$.
\end{proposition}

\begin{proof}
This is the propositional restatement of Theorem~\ref{thm:edge-34}
used by citation sites in Chapter~\ref{ch:chiral-higher-deligne} and
Chapter~\ref{ch:brace}; the content is identical. The factorization
algebra $\Obs^{\mathrm{q}}$ of the bulk--boundary BV theory restricts
on the closed half to the factorization face $\Obs^{\mathrm{ch}}_\cA$
(Theorem~\ref{thm:edge-34}), and on a cofibrant replacement the QME
equals the convolution MC (Theorem~\ref{thm:edge-45}); composing these
two quasi-isomorphisms with the coloured Swiss-cheese formality map
(Getzler--Jones~1994, coloured variant via Vallette~2014 colored Koszul
duality plus homotopy transfer) produces the stated quasi-isomorphism
between the factorization brace action and the convolution
$L_\infty$-algebra. Strictness on a fixed semifree model and
associator-dependence follow from Theorem~\ref{thm:edge-45} and the
$\mathrm{GRT}_1$-torsor structure of Swiss-cheese formality.
\end{proof}

\subsection{Edge (4)--(5): BV/BRST--convolution}\label{subsec:edge-45}

\begin{theorem}[QME equals convolution MC; \ClaimStatusProvedHere]
\label{thm:edge-45}%
\index{heptagon!edge (4)--(5)}
On a fixed semifree model (i.e.\ cofibrant replacement of the
BV/BRST complex), the QME of Face~(4) is strictly equal to the
Maurer--Cartan equation in $\fg^{\,\SCchtop}_{\cA,Z}$ of Face~(5).
Passing to quasi-iso classes recovers the $\infty$-quasi-isomorphism
of FM160.
\end{theorem}

\begin{proof}
Fix a semifree model $\cE\to\SCchtop$ of coloured $\infty$-operads
and a coordinate system on the BV complex. The QME in those
coordinates is a polynomial in the structure constants of
$\SCchtop$ acting on the BV field content, which is, by
construction of the convolution algebra, the MC equation in
$\fg^{\,\SCchtop}_{\cA,Z}=\mathrm{Hom}^{\mathrm{col}}(B(\SCchtop),
\End^{\mathrm{col}}_{(\cA,Z)})$. The equality is strict on the
semifree model; on arbitrary models the two sides are related by
homotopy transfer (HPL). Gauge-fixed chain-level identity closes
FM160.
\end{proof}

\begin{proposition}[Heptagon edge $(4)\leftrightarrow(5)$, propositional
form; \ClaimStatusProvedHere]\label{prop:heptagon-edge-45}%
\index{heptagon!edge (4)--(5) propositional form}%
The BV/BRST presentation (Face~(4)) and the convolution
$L_\infty$-presentation (Face~(5)) of $\SCchtop$ yield the same
$\infty$-operadic datum on $(\cA, Z^{\mathrm{der}}_{\mathrm{ch}}(\cA))$.
Concretely: the quantum master equation on the observables
$\Obs^{\mathrm{q}}(U)$ descends to the Maurer--Cartan equation in
$\fg^{\,\SCchtop}_{\cA,Z}$ via a gauge-fixed semifree model; the
quasi-isomorphism class agreement is proved by Stasheff cancellation
on the evaluation-generated core of the coloured operad.
Consequently, the Dunn-additive product
$E_2\text{-chiral}\otimes_{\mathrm{Dunn}} E_1^{\mathrm{top}}$
identifies, through Face~(4)--Face~(5), with a single $E_3$-action
on $Z^{\mathrm{der}}_{\mathrm{ch}}(\cA)$.
\end{proposition}

\begin{proof}
The first sentence is the propositional restatement of
Theorem~\ref{thm:edge-45}; content is identical. For the Dunn claim:
the convolution algebra $\fg^{\,\SCchtop}_{\cA,Z}$ carries the
$E_2$-chiral structure from the closed-colour braces
(Theorem~\ref{thm:edge-23}) and an independent $E_1$-topological
structure from the conformal-weight $L_0$-flow along $\R$; Dunn
additivity for chain-level factorisation algebras (Costello--Gwilliam
Vol.~II~\S5; Lurie HA~5.1.2.2 applied to coloured operads with
empty cross-arrows at the mixed sector) combines them into the
$E_3$-action. Associator-dependence and Stasheff-cancellation on the
evaluation-generated core (Vol.~I, AP47) produce the quasi-iso class
equality on arbitrary models.
\end{proof}

\subsection{Edge (5)--(6): convolution--Drinfeld centre}\label{subsec:edge-56}

\begin{theorem}[Convolution MC presents categorical half-braiding;
\ClaimStatusProvedHere]\label{thm:edge-56}%
\index{heptagon!edge (5)--(6)}
The Maurer--Cartan locus of the convolution $L_\infty$-algebra
$\fg^{\,\SCchtop}_{\cA,Z}$ is in natural bijection with the space of
half-braidings $\sigma_\cA(z)$ compatible with the
$\SCchtop$-datum on $(\cA,Z^{\mathrm{der}}_{\mathrm{ch}}(\cA))$, hence
with the central objects of
$\cZ(\Rep_{\mathrm{fact}}(\cA))$ of Face~(6):
\begin{equation}
  \MC\bigl(\fg^{\,\SCchtop}_{\cA,Z}\bigr)
  \;\simeq\;
  \bigl(\mathrm{objects\ of\ }
  \cZ(\Rep_{\mathrm{fact}}(\cA))\bigr).
\end{equation}
Gauge equivalence in convolution corresponds to isomorphism in
the Drinfeld centre.
\end{theorem}

\begin{proof}
A Maurer--Cartan element $\alpha\in\fg^{\,\SCchtop}_{\cA,Z}$
determines a coherent family of operations of $\SCchtop$ on
$(\cA,Z)$; the mixed-sector component of $\alpha$ is precisely a
half-braiding (by Proposition~\ref{prop:half-braiding-construction};
the mixed sector of $\SCchtop$'s Koszul dual is the shuffle-mixed
data, which translates to half-braidings via
Definition~\ref{def:chiral-half-braiding}). Conversely, a central
object of $\cZ(\Rep_{\mathrm{fact}}(\cA))$ gives a half-braiding
compatible with $\Delta_z$, which encodes an MC element of the
mixed sector of $\fg^{\,\SCchtop}_{\cA,Z}$. Gauge equivalence in
convolution is homotopy of MC elements
(Hinich~1997, Getzler~2009); this corresponds to isomorphism in
$\cZ(\Rep_{\mathrm{fact}}(\cA))$ by Theorem~\ref{thm:drinfeld-centre-sc-face}.
\end{proof}

\subsection{Edge (6)--(7): Drinfeld centre--derived algebraic geometry}\label{subsec:edge-67}

\begin{theorem}[Drinfeld centre is the
$\infty$-stack of centrally-compatible $\SCchtop$-data;
\ClaimStatusProvedHere]\label{thm:edge-67}%
\index{heptagon!edge (6)--(7)}
The Drinfeld centre $\cZ(\Rep_{\mathrm{fact}}(\cA))$ of Face~(6)
is the loop space of $\cM_{\SCchtop,X}$ at the point
$(\cA,Z^{\mathrm{der}}_{\mathrm{ch}}(\cA))$, in the sense of
$E_2$-group objects:
\begin{equation}
  \label{eq:edge-67-equivalence}
  \cZ\bigl(\Rep_{\mathrm{fact}}(\cA)\bigr)
  \;\simeq\;
  \Omega_{(\cA,Z)}^{E_2}\,\cM_{\SCchtop,X}
\end{equation}
as braided monoidal $\infty$-categories. The $1$-shifted symplectic
form $\omega_{\SCchtop}$ of Face~(7) restricts on the loop space to
the categorical braiding of Face~(6).
\end{theorem}

\begin{proof}
The mapping-stack tangent $T_{(\cA,Z)}\cM_{\SCchtop,X}$ is
$\fg^{\,\SCchtop}_{\cA,Z}$ by
Theorem~\ref{thm:convolution-lagrangian}. Its $E_2$-group
completion (iterated based loop space) is, by Deligne conjecture
for coloured operads (coloured chiral higher Deligne;
Theorem~\ref{thm:chiral-higher-deligne} of
Chapter~\ref{ch:theory-hochschild} applied to $(\SCchtop)^!$),
the $E_2$-algebra whose $\Rep$ category is
$\cZ(\Rep_{\mathrm{fact}}(\cA))$. This gives the equivalence
in~\eqref{eq:edge-67-equivalence}.

The restriction of $\omega_{\SCchtop}$ to the loop space:
PTVV Corollary~2.6 states that for a $k$-shifted symplectic derived
stack $\cX$, the loop space $\Omega\cX$ is $(k-1)$-shifted
symplectic; applied to $k=1$, $\Omega\cM_{\SCchtop,X}$ is
$0$-shifted symplectic, i.e.\ carries a symmetric non-degenerate
pairing in degree $0$, which is exactly the categorical braiding
pairing (symmetric at degree $0$, producing the full hexagon at
all orders via Deligne's coherence). This identifies the braiding
of Face~(6) with the loop-space shifted symplectic form.
\end{proof}

\subsection{Edge (7)--(1): derived algebraic geometry--operadic}\label{subsec:edge-71}

\begin{theorem}[Tangent complex of $\cM_{\SCchtop,X}$ recovers the
operadic face; \ClaimStatusProvedHere]\label{thm:edge-71}%
\index{heptagon!edge (7)--(1)}
At every closed point $(\cA,Z)\in\cM_{\SCchtop,X}$, the
tangent dg-Lie algebra equipped with its coloured operad structure
is a strict model of $\SCchtop$ on the pair
$(\cA,Z^{\mathrm{der}}_{\mathrm{ch}}(\cA))$. The global
sections of the structure sheaf
$\mathcal{O}_{\cM_{\SCchtop,X}}$ recover the coloured cobar
construction $\Omega(B(\SCchtop))$ on any local model, hence present
Face~(1) directly.
\end{theorem}

\begin{proof}
Tangent dg-Lie identification is
Theorem~\ref{thm:convolution-lagrangian}. The coloured $\SCchtop$-operad
structure on the tangent complex is inherited from the operadic
ambient of Definition~\ref{def:moduli-sc-ch-top}: the mapping-stack
target $B\SCchtop\mathrm{-Alg}$ has tangent $\SCchtop$ itself (up to
homotopy; Lurie HA~5.3.1.11). Pulling back along the tautological
evaluation map, the tangent of $\cM_{\SCchtop,X}$ carries
$\SCchtop$ up to $\infty$-quasi-isomorphism.

Global sections of $\mathcal{O}_{\cM_{\SCchtop,X}}$: by the
Kan extension of the tautological family to the moduli stack,
the Chevalley--Eilenberg complex of the tangent dg-Lie is
$\mathcal{O}^{\mathrm{loc}}_{\cM_{\SCchtop,X}}$ (locally
around any closed point; PTVV 2013~\S2.2 and Deligne's
\textit{Equivalence of Deformation Complexes}). The
CE complex is, by construction, $\Omega(B(\SCchtop))=W(\SCchtop)$
applied to the tangent, hence recovers Face~(1) via the edge
(1)--(2) Koszul-dual equivalence.
\end{proof}

%%%%%%%%%%%%%%%%%%%%%%%%%%%%%%%%%%%%%%%%%%%%%%%%%%%%%%%%%%%%%%%%%%%%%%%%%%%%%
\section{Closing the heptagon}
\label{sec:heptagon-closed}
%%%%%%%%%%%%%%%%%%%%%%%%%%%%%%%%%%%%%%%%%%%%%%%%%%%%%%%%%%%%%%%%%%%%%%%%%%%%%

\begin{theorem}[Heptagon of $\SCchtop$; \ClaimStatusProvedHere]
\label{thm:heptagon-closed}%
\index{heptagon!seven-way equivalence|textbf}
The seven faces (1)--(7) of $\SCchtop$ are pairwise
$\infty$-quasi-isomorphic as coloured $\infty$-operads after
fixing a Drinfeld associator $\Phi\in\mathrm{GRT}_1(\mathbb{Q})$,
with edges given by Theorems~\ref{thm:edge-12}, \ref{thm:edge-23},
\ref{thm:edge-34}, \ref{thm:edge-45}, \ref{thm:edge-56},
\ref{thm:edge-67}, \ref{thm:edge-71}. The seven-way equivalence is
$\Phi$-independent at the level of cohomology; strictly chain-level
presentations track a single $\Phi$.
\end{theorem}

\begin{proof}
Compose the seven edges around the heptagon in either direction.
Each edge is an $\infty$-quasi-isomorphism
(Convention~\ref{conv:edge-regime}); composition preserves
$\infty$-quasi-isomorphisms (Fresse 2009, or direct verification in
the colored $\infty$-operad model category); the composed
self-equivalence of Face~(1) is homotopic to the identity because
each edge is canonical modulo associator choice. $\Phi$-independence
at cohomology follows from the GRT$_1$-action on each edge
(Willwacher 2015; coloured extension via coloured brace
formality).
\end{proof}

\begin{remark}[The fourteen diagonals]\label{rem:heptagon-diagonals}
The fourteen diagonal edges of the heptagon are compositions of
the seven adjacent edges; each records a direct equivalence between
non-adjacent faces. Notable diagonals include: Face~(1)--(3) (operadic
action on factorization algebra via chiral Deligne--Tamarkin;
Chapter~\ref{ch:theory-hochschild}); Face~(2)--(5) (Koszul dual
as tangent dg-Lie; Vallette Koszul duality); Face~(3)--(7)
(factorization algebra as global sections of $\cM_{\SCchtop,X}$);
Face~(1)--(6) (direct bar--cobar to Drinfeld centre via
Theorem~\ref{thm:drinfeld-centre-sc-face} after composing with
Edge~(1)--(2)); Face~(4)--(7) (QME as defining equation of a
Lagrangian section of the shifted symplectic stack). Each diagonal
has an independent direct construction in the literature; the
heptagon gathers them under one roof.
\end{remark}

%%%%%%%%%%%%%%%%%%%%%%%%%%%%%%%%%%%%%%%%%%%%%%%%%%%%%%%%%%%%%%%%%%%%%%%%%%%%%
\section{Conformal vector and the collapse to $E_3$-topological}
\label{sec:heptagon-conformal-vector}
%%%%%%%%%%%%%%%%%%%%%%%%%%%%%%%%%%%%%%%%%%%%%%%%%%%%%%%%%%%%%%%%%%%%%%%%%%%%%

The heptagon is the \emph{generic} picture. Adding a conformal vector
to the boundary collapses two faces of the heptagon into one; the
resulting topologised structure is $E_3$-topological, not $\SCchtop$
(Chapter~\ref{ch:topologization}; AP-TOPOLOGIZATION).

\begin{theorem}[Collapse with conformal vector at affine Kac--Moody
non-critical level; \ClaimStatusProvedHere]
\label{thm:heptagon-collapse}%
\index{topologisation!collapses heptagon|textbf}%
\index{$E_3$-topological!heptagon collapse}
Let $\cA=V_k(\fg)$ be an affine Kac--Moody chiral algebra at
non-critical level $k\neq -h^{\vee}$, with Sugawara conformal
vector $T(z)=\tfrac{1}{2(k+h^{\vee})}:\!J^aJ_a\!:(z)$. In the
cohomology of the BRST differential $Q_{\mathrm{tot}}$
(Chapter~\ref{ch:topologization}), the two colours of $\SCchtop$
merge: $T(z)$ satisfies
$T(z)=[Q_{\mathrm{tot}},G(z)]$ for a BRST primitive $G(z)$, making
$C$-translations $Q_{\mathrm{tot}}$-exact; the closed-colour
$\FM_k(\mathbb{C})$ and the open-colour $\Conf_m(\mathbb{R})$
collapse into a single $\FM_{k+m}(\mathbb{R}^3)$ via
Dunn additivity $E_2^{\mathrm{hol}}\otimes E_1^{\mathrm{top}}\simeq
E_3^{\mathrm{top}}$ in cohomology. The heptagon's seven faces
collapse to:
\begin{enumerate}
  \item operadic face for $E_3$-topological $\FM(\mathbb{R}^3)$;
  \item Koszul-dual face $E_3^{!}\simeq E_3\{3\}$;
  \item factorization face on $X\times\mathbb{R}\simeq\mathbb{R}^3$
  (after choosing a point on $X$);
  \item BV/BRST face of the resulting $3$-dimensional topological
  field theory;
  \item convolution $L_\infty$-algebra for $E_3^{!}$-algebras;
  \item Drinfeld centre of a module $1$-category over the
  mode algebra of $\cA$;
  \item derived stack of $3$-dimensional topological algebras.
\end{enumerate}
The seven $E_3$-topological faces remain quasi-isomorphic under the
topologisation functor (AP-TOPOLOGIZATION).
\end{theorem}

\begin{proof}
Sugawara's identity writes $T(z)$ as a quadratic Casimir; the BRST
primitive $G(z)$ is constructed explicitly from the antighost
bilinear (Chapter~\ref{ch:topologization}
Theorem~\ref{thm:E3-topological-km}; Costello 2016;
CFG~\cite{CFG2602}). The resulting $Q_{\mathrm{tot}}$-cohomology
kills the holomorphic direction; Dunn additivity of little discs
(Lurie HA~5.1.2.2, Dunn 1988) passes
$E_2^{\mathrm{hol}}\otimes E_1^{\mathrm{top}}$ to $E_3^{\mathrm{top}}$
at the level of operads. Each heptagon face inherits the
topologisation via its edge to Face~(3) or Face~(1); the resulting
$E_3$-topological faces are equivalent among themselves by the
standard $E_3$-operad Koszul duality (Fresse Vol.~II).
\end{proof}

\begin{conjecture}[Collapse with conformal vector, general case]
\label{conj:heptagon-collapse-general}%
\index{topologisation!general chiral with conformal|textbf}
Let $\cA$ be an $E_1$-chiral algebra carrying a conformal vector
$T(z)$ at non-critical level, not necessarily of affine Kac--Moody
type. Then the heptagon collapses to an $E_3$-topological
seven-face structure as in Theorem~\ref{thm:heptagon-collapse}.
Evidence: the construction of the BRST primitive $G(z)$ is
conjectural (AP-TOPOLOGIZATION); the chain-level Dunn additivity
is conjectural; the class-$M$ chain-level case
(Virasoro, $W_N$) is the principal open direction of
Chapter~\ref{ch:topologization}. See
Conjecture~\ref{conj:topologization-general}.
\end{conjecture}

\begin{remark}[The heptagon survives the generic case]
\label{rem:heptagon-survives-generic}
For chiral algebras without a conformal vector (class $E_1$-chiral
Yangians, critical-level Kac--Moody, CY functor outputs without
conformal vector), all seven faces of the heptagon remain distinct
and the heptagon is the final structure. This is the generic case;
$E_3$-topologisation is the exceptional special case triggered by a
conformal vector at non-critical level. The programme's climax is
the heptagon, not the topologised special case.
\end{remark}

\begin{remark}[Critical level]\label{rem:heptagon-critical}
At critical level $k=-h^{\vee}$ for affine Kac--Moody, the Sugawara
tensor $T(z)$ degenerates: the normalisation factor
$(k+h^{\vee})^{-1}$ diverges; the BRST primitive $G(z)$ does not
exist. The heptagon stays intact, but the boundary chiral algebra
is enriched by the Feigin--Frenkel centre of
$\hat{\mathfrak{g}}_{-h^{\vee}}$; the fifth shadow class FF
(Feigin--Frenkel) of Infinite Fingerprint Classification emerges
(Platonic Reconstitution). Topologisation fails; the heptagon is
the complete structural statement.
\end{remark}

%%%%%%%%%%%%%%%%%%%%%%%%%%%%%%%%%%%%%%%%%%%%%%%%%%%%%%%%%%%%%%%%%%%%%%%%%%%%%
\section{Summary table and commitments}
\label{sec:heptagon-summary}
%%%%%%%%%%%%%%%%%%%%%%%%%%%%%%%%%%%%%%%%%%%%%%%%%%%%%%%%%%%%%%%%%%%%%%%%%%%%%

\begin{table}[h]
\centering
\begin{tabular}{rllc}
\toprule
Face & Presentation & Home category & Status \\
\midrule
1 & Operadic: $C_\bullet(\FM\times\Conf)$ & Coloured homotopy operads
& ProvedHere (\ref{subsec:face-operadic}) \\
2 & Koszul-dual: $(\Lie,\Ass,\text{shuffle-mix})$ & Coloured Koszul
operads & ProvedHere (\ref{subsec:face-koszul-dual}) \\
3 & Factorization: $\Obs^{\mathrm{ch}}$ + brace & $\Dmod(\Ran X)$ &
ProvedHere (\ref{subsec:face-factorization}) \\
4 & BV/BRST: $\Obs^q$ on $X\times\mathbb{R}_{\geq 0}$ &
Log $\SCchtop$-algebras & ProvedHere (\ref{subsec:face-bv-brst}) \\
5 & Convolution: $\fg^{\,\SCchtop}_{\cA,Z}$ & Coloured
$L_\infty$-algebras & ProvedHere (\ref{subsec:face-convolution}) \\
6 & Drinfeld centre: $\cZ(\Rep_{\mathrm{fact}}(\cA))$ &
$E_2$-$\mathrm{Cat}_\infty$ & ProvedHere
(Thm.~\ref{thm:drinfeld-centre-sc-face}) \\
7 & Derived stack: $\cM_{\SCchtop,X}$ & $1$-shifted
$\mathrm{Sp}\,\infty$-stacks & ProvedHere
(Thms.~\ref{thm:moduli-sc-shifted-symplectic},
\ref{thm:convolution-lagrangian}) \\
\bottomrule
\end{tabular}
\caption{The seven faces of the heptagon. Every face presents the
same coloured $\infty$-operad $\SCchtop$; adjacent faces are joined
by the seven edges of \S\ref{sec:heptagon-edges}; the closure is
Theorem~\ref{thm:heptagon-closed}.}
\label{tab:heptagon-faces}
\end{table}

\begin{remark}[Associator-free reading]\label{rem:heptagon-assoc-free}
The cohomology class of the heptagon is associator-free. A single
choice of Drinfeld associator $\Phi\in\mathrm{GRT}_1(\mathbb{Q})$
fixes a strict chain-level presentation; varying $\Phi$ transports
the presentation through the $\mathrm{GRT}_1$-torsor action (FM158
heal). The Koszulness Moduli atlas (Platonic Reconstitution,
Chapter~\ref{ch:koszulness-moduli}) assigns one chart per
characterisation; Face~$i$'s natural coordinate associator is the
Drinfeld--Tamarkin associator canonical to that presentation.
\end{remark}

\begin{remark}[Beyond the heptagon]\label{rem:heptagon-beyond}
The heptagon is not a final horizon. Two further faces are natural
candidates: Face~(8), the motivic face via Brown's motivic coaction
(GRT-parametrized Seven Faces, Platonic Reconstitution); Face~(9),
the operadic graph-complex face via Willwacher's
$H^0(\mathrm{GC}_2)=\mathfrak{grt}_1$ identification. These form
the head of an infinite $\mathrm{GRT}_1$-torsor of faces. The
heptagon is a rigidly-determined coarse slice; the full torsor is
the subject of Chapter~\ref{ch:grt-faces}.
\end{remark}

%%%%%%%%%%%%%%%%%%%%%%%%%%%%%%%%%%%%%%%%%%%%%%%%%%%%%%%%%%%%%%%%%%%%%%%%%%%%%
\section*{Notes and provenance}
%%%%%%%%%%%%%%%%%%%%%%%%%%%%%%%%%%%%%%%%%%%%%%%%%%%%%%%%%%%%%%%%%%%%%%%%%%%%%

Face~(1) originates with Voronov's Swiss-cheese operad (1998); the
chiral upgrade is Beilinson--Drinfeld \textit{Chiral Algebras}~\S3.
Face~(2) is Hoefel~2009 and Hoefel--Livernet~2012; the three-sector
structure is established in Chapter~\ref{ch:theory-equivalence}.
Face~(3) is Costello--Gwilliam Vol.~II; brace structure is Tamarkin
2000. Face~(4) is Costello--Gwilliam Vol.~I and Butson--Yoo 2018;
logarithmic refinement is Chapter~\ref{ch:bv-construction}. Face~(5)
is Vallette coloured Koszul duality. Face~(6) is new to this chapter,
built on the half-braiding construction of
Definition~\ref{def:half-braiding-explicit} of
Chapter~\ref{ch:theory-hochschild}. Face~(7) is
PTVV~2013 applied to the mapping stack of
Definition~\ref{def:moduli-sc-ch-top}.

The seven-fold equivalence is new; partial equivalences were known
classically (edges~(1)--(2), (1)--(3), (3)--(4), (4)--(5)). The
upgrade from pentagon to heptagon is the present chapter's
contribution.
